\makeabstract
    {
    Effects of Repeat Wildfires on Boreal Forest Seedlings
    }
    {
    Cooper Merrill\textsuperscript{1}, Austin McGuire\textsuperscript{1}, Wylee Drew\textsuperscript{1}, Miranda Yang\textsuperscript{1}, Katie Nelson\textsuperscript{1}, Rebecca Hewitt\textsuperscript{1}, Fleur Nicklen\textsuperscript{2}, Paige Paulsen\textsuperscript{2}
    }
    {
    \textsuperscript{1} Department of Environmental Studies, Amherst College, Amherst, MA\\
\textsuperscript{2} Eastern Area Fire Management, National Parks Service, Fairbanks, AK
    }
    {
    Boreal forests in interior Alaska are burning more frequently as the climate warms, raising questions about the future of  ecosystems whose dominant species are adapted to historically long fire-return intervals. Because ectomycorrhizal fungi  associated with fine roots are an important part of nutrient acquisition for Black Spruce (Picea mariana) and other boreal forest species, we asked how repeated fire reshapes plant-fungal nutrient acquisition strategies. 
    }
    {
    Within Wrangell St. Elias National Park, we sampled plots that had burned once (in 2009) or twice (2009 and 2016). From each plot, we collected soil samples and root and foliar samples from seedlings of the four dominant tree species: Black Spruce, White Spruce (Picea glauca), Trembling Aspen (Populus tremuloides), and Alaskan Paper Birch (Betula neoalaskana). Seedling traits were recorded, and carbon and nitrogen percentages and isotope compositions were analyzed in soil and foliage samples. We analyzed roots for morphological traits such as length, diameter, volume, number of root tips, and number of forks using the WinRHIZO software and picked root tips colonized by ectomycorrhizal fungi for DNA extraction.
    }
    {
    Linear mixed-effects models showed that, in twice-burned plots, trembling aspen produced longer, thinner roots with more root tips, consistent with a shift toward a more independent, less mycorrhizal dependent nutrient-acquisition strategy in response to the shortened burn interval. Root morphology did not differ significantly with burn history in White Spruce, Black Spruce, and Alaskan Paper Birch; however, low seedling recruitment (n \ensuremath{<} 10 for some sample groups) limited statistical power. 
    }
    {
    Together these results suggest that shortened wildfire intervals could cause Trembling Aspen to reduce their investment in mycorrhizal fungi. Analysis of nutrient and isotope data  as well as Molecular identification of the fungi colonizing root tips is underway and will allow us to test how the belowground community responds to shortened wildfire intervals.
    }

    \newpage
    